\section{Введение}
\label{sec:intro}

Эта статья описывает одну машинно-проверенную идею и дисциплину, которая держит её честной.
Идея --- физический запрет: двигатель, который \emph{вечно} совершает нетривиальную работу
спуска, ничего при этом не теряя, существовать не может. В арифметике натуральных чисел это
метод бесконечного спуска Ферма; мы переформулируем его для центров $m$ евклидовых пар
$(6m-1,\,6m+1)$ и называем запрещённый объект \emph{вечным двигателем Евклида}. Сам запрет ---
не существует бесконечно убывающей цепи чистых евклидовых шагов --- доказан на голом ядре
Lean~4, без \code{mathlib} и без каких-либо аксиом сверх стандартных (\cref{sec:engine}).

Тезис программы состоит в том, что удивительно многие трудные вопросы суть тени этого одного
запрета, отброшенные на разные объекты. Там, где объект не может нести вечный двигатель,
структурная теорема «нет отклонения» выполняется \emph{безусловно}; открытым же остаётся всегда
одно и то же --- связь между абстрактной структурой и конкретным аналитическим объектом,
которого математика пока не предоставила. Мы проводим это прочтение через гипотезу о
простых-близнецах (стержень статьи), гипотезу Римана, $\mathrm{P}$ против $\mathrm{NP}$,
Янга--Миллса, Ходжа, Навье--Стокса, Бёрча--Свиннертон-Дайера и далее --- через простые Мерсенна,
арифметический зоопарк узоров из простых, геометрию графа спуска и задачу Коллатца.

\paragraph{Что доказано, а что нет.}
Мы настойчивы в этом пункте, потому что тема провоцирует преувеличения, а мы не хотим ни
одного. Гипотеза о простых-близнецах здесь \emph{не} доказана. Доказана же, и машинно, без
\code{sorry}, \emph{редукция}: вся близнецовая ветвь сведена к одному явному открытому узлу ---
конечному семантическому ключу, который должен разрешать определённые коллизии генеалогий
(\cref{sec:reduction}). Этот узел нельзя закрыть изнутри системы: его закрытие означало бы запуск
того самого двигателя, невозможность которого лежит в основании. Поэтому мы принимаем его
\emph{извне}, как единственную аксиому, называемую \emph{первопричиной} (\cref{sec:firstcause}),
запертую в одном карантинном модуле. Наша главная теорема, \emph{высшая энергетическая
несовместимость}, делает картину точной: бесконечность близнецов нельзя ни доказать, ни
опровергнуть внутри системы без вечного двигателя, так что конечный наблюдатель может лишь
\emph{принять} первопричину, но никогда её не вывести.

\paragraph{Дисциплина честности.}
Через всё построение проходят три цвета, и мы называем их прямо: \green{}~утверждение
машинно-доказано при стандартных аксиомах (сюда входит и \emph{зелёная условная теорема}, чьё
условие --- названный вход, а не аксиома); \yellow{}~утверждение \emph{заражено аксиомой}, то
есть зависит от первопричины; \red{}~объект действительно открыт. Отдельный верификатор обходит
каждую декларацию репозитория, собирает аксиомы, от которых она зависит, и обеспечивает два
неизменных инварианта: единственная нестандартная аксиома во всём построении --- \axiom{}, а
единственное заражённое \code{sorry} утверждение --- сама гипотеза о близнецах. На момент письма
верификатор сообщает ровно одну нестандартную аксиому, ровно один \code{sorry} и $47$
заражённых аксиомой деклараций; зелёный модуль подделать нельзя, а аксиом-трекер не оставляет
декрету места спрятаться за зелёной вывеской. Четырежды в ходе работы внутренний аудит находил
пустую обёртку раньше, чем она попадала в витрину (мы отмечаем эти вакуумности там, где они
возникли); то, что они пойманы машиной, --- лучший довод доверять частям, прошедшим витрину.

\paragraph{Воспроизводимость.}
Всё в этой статье --- отражение исходного кода Lean~4~\cite{lean4,mathlib,euclidspath-repo}.
Каждое формальное утверждение ниже помечено именем своей машинно-проверенной декларации,
набранным \code{машинописью}. Полный исходный код открыт по адресу
\url{https://github.com/elamaunt/Euclids-path}; идентичный контент доступен как двуязычный сайт
по адресу \url{https://elamaunt.github.io/Euclids-path/}. Сборка командой \code{lake build}
проверяет всё построение, а \code{lake env lean scripts/VerifyAll.lean} заново устанавливает два
инварианта честности и громко падает, если нарушен любой из них.

\paragraph{План.}
\cref{sec:engine} доказывает невозможность двигателя и его необратимость. \cref{sec:carrier}
фиксирует носитель двойки и законы сохранения, делающие пары «разделёнными». \cref{sec:reduction}
сводит гипотезу о близнецах к одному узлу. \cref{sec:firstcause} вводит первопричину и доказывает
главную теорему. \cref{sec:riemann,sec:pnp,sec:ymhodge,sec:ns} прочитывают гипотезу Римана,
$\mathrm{P}$ против $\mathrm{NP}$, Янга--Миллса и Ходжа, а также Навье--Стокса через двигатель.
\cref{sec:zoo,sec:geometry,sec:bsd} разбирают ветвь Мерсенна и арифметический зоопарк, геометрию
пути и Бёрча--Свиннертон-Дайера вместе с несколькими открытыми соседями. \cref{sec:collatz}
разбирает задачу Коллатца, включая декрет, который был взят и затем снят, когда его закон был
машинно опровергнут. \cref{sec:coda} --- синтез; \cref{sec:status} даёт честный глобальный статус
и процедуру верификации; \cref{sec:figures} задаёт алгоритм генерации каждого рисунка.
